%% @Author: Ahmad Ben Taleb
%  @Date:   2026
%% @Class:  PFE ITBS V3.

\documentclass[a4paper, oneside, 12pt, final]{extreport}
\usepackage{graphicx}
\usepackage{float}
\parindent 0cm
\usepackage{makeidx}
\usepackage{tcolorbox}
\makeindex
\usepackage[lined,boxed,commentsnumbered, english, ruled,vlined,linesnumbered]{algorithm2e}
\usepackage{amsthm}
\newtheorem{theorem}{Theorem}[chapter]
\newtheorem{definition}{Definition}[chapter]
\newtheorem{exemple}{Example}[chapter]
\usepackage[nottoc]{tocbibind}
\providecommand{\keywords}[1]{\textbf{\textit{Mots clés---}} #1}
\providecommand{\keywordss}[1]{\textbf{\textit{Keywords---}} #1}
\usepackage{etoolbox}
\textwidth 18cm
\textheight 24cm
\topmargin -0.5cm
\oddsidemargin -1cm

\usepackage{ifxetex}
\ifxetex
\usepackage{fontspec}
\else
\usepackage[T1]{fontenc}
\usepackage[utf8]{inputenc}
\usepackage{lmodern}
\fi

% ──────────────────────────────────────────────────────────
%  PAGE DE GARDE
% ──────────────────────────────────────────────────────────
\newcommand{\reportTitle}{%
\textsc{Optical Character Recognition: Deep Learning Approach for Printed Text}}

\newcommand{\reportAuthor}{%
Chebâane  \textsc{Mohamed Yassine} \\
Mbarek \textsc{Issam idin}%
}

\newcommand{\reportSubject}{}

\newcommand{\studyDepartment}{%
Entreprise d'accueil}

\newcommand{\ITBS}{%
\\ Ecole Supérieure Privée des Technologies de l'Information et de Management de Nabeul}

\newcommand{\AU}{\centering \textbf{Année Universitaire 2024-2025}}

\newcommand{\specialcell}[1]{%
\begin{tabularx}{\textwidth}{@{}X@{}}#1\end{tabularx}%
}

\newcommand{\MyCommand}{Does nothing really}

\title{\reportSubject}
\author{\reportAuthor}

\usepackage{graphics}
\usepackage{graphicx}
\usepackage[acronym,toc,section=chapter]{glossaries}
\makeglossaries

\newacronym{ocr}{OCR}{Optical Character Recognition}
\newacronym{cnn}{CNN}{Convolutional Neural Network}
\newacronym{rnn}{RNN}{Recurrent Neural Network}
\newacronym{lstm}{LSTM}{Long Short-Term Memory}
\newacronym{ctc}{CTC}{Connectionist Temporal Classification}
\newacronym{crnn}{CRNN}{Convolutional Recurrent Neural Network}
\newacronym{bilstm}{BiLSTM}{Bidirectional Long Short-Term Memory}
\newacronym{svm}{SVM}{Support Vector Machines}
\newacronym{hdf5}{HDF5}{Hierarchical Data Format version 5}
\newacronym{api}{API}{Application Programming Interface}

\pagenumbering{roman}
\usepackage[utf8]{inputenc}

\begin{document}

\thispagestyle{empty}

\begin{titlepage}
\begin{center}

%% ── HEADER ──────────────────────────────────
\includegraphics[scale=0.20]{embleme.jpg}
\vspace{0.5cm}
{%
\fontsize{9pt}{9pt}\selectfont%
\begin{tabular}{c}
R\'epublique Tunisienne \\
Minist\`ere de l'Enseignement Supérieur et de la Recherche Scientifique   
\ITBS{}
\\ 
\end{tabular}
}

\vspace{1cm}
\includegraphics[scale=0.1]{logonoir.png}

%% ── CONTENT ────────────────────────────────────────────
\vspace{10pt}
\vspace{5pt}
\vspace{15pt}
{\textit{Mini projet Deep Learning}}\\
\vspace{10pt}
{\textbf{\large Diplôme National d'Ingénieur en Génie Informatique \\ Génie logiciel}}\\
\vspace{5pt}

\textbf{\textit{Réalisé par}}\\
\vspace{10pt}
{%
\fontsize{14pt}{14pt}\selectfont%
{\bfseries\Large\sc \reportAuthor}\\
}%

\vspace{5pt}
{%
\fontsize{27pt}{27pt}\selectfont%
\rule{0.5\textwidth}{.4pt}\\
\vspace{10pt}
\reportSubject{Optical Character Recognition (OCR):\\
A Deep Learning Approach for Printed Text Recognition}\\
%
\vspace{10pt}
\rule{0.5\textwidth}{.4pt}
}

\vspace{5pt}
\vspace{10pt}

\begin{table}[h]
\begin{tabular}{lcr}
\textbf{Encadrant Académique:} M Ben Taleb Ahmed      & 
\end{tabular}
\end{table}

\vspace{10pt}
\textbf{\textit{}}\\
\vspace{5pt}
\studyDepartment

\centering
\includegraphics[scale=0.08]{logonoir.png}

\end{center}
\vspace{20pt}
\AU\\

\end{titlepage}

% ──────────────────────────────────────────────────────────
%  DÉDICACE
% ──────────────────────────────────────────────────────────
\chapter*{Dedication}
\thispagestyle{empty}

\begin{center}
{\it
We dedicate this work to our families, for their unconditional support, encouragement,
and patience throughout our studies. Their guidance and trust have been a constant
source of motivation.

We also dedicate this project to our friends, who have always been there to inspire us,
share their ideas, and encourage us along the way\ldots
}
\end{center}

% ──────────────────────────────────────────────────────────
%  REMERCIEMENTS
% ──────────────────────────────────────────────────────────
\chapter*{Acknowledgment}
\thispagestyle{empty}

We would like to express our sincere gratitude to all those who have supported us
throughout this project.

First and foremost, we would like to thank our academic supervisor, Mr.\ Ben Taleb Ahmed,
for his guidance, valuable advice, and continuous encouragement. His expertise and
patience have been instrumental in the successful completion of this work.

We are also grateful to our colleagues and friends who offered their support, shared ideas,
and provided constructive feedback during the development of this project.

Finally, we wish to thank our families for their unwavering support, understanding, and
encouragement throughout our studies. Their trust and motivation have been a constant
source of strength.

This work would not have been possible without all of their contributions.

% ──────────────────────────────────────────────────────────
%  TABLE DES MATIÈRES
% ──────────────────────────────────────────────────────────
\addtocontents{toc}{\protect\setcounter{tocdepth}{4}}
\tableofcontents
\listoffigures
\listoftables
\printglossaries

\cleardoublepage
\newpage
\pagenumbering{arabic}

% ──────────────────────────────────────────────────────────
%  INTRODUCTION GÉNÉRALE
% ──────────────────────────────────────────────────────────
\chapter*{Introduction}
\label{chap:general_introduction}
\addcontentsline{toc}{chapter}{Introduction}
\section*{Context and Motivation}

In the digital age, the ability to automatically extract text from images has become
increasingly important across numerous domains. \gls{ocr} technology enables the conversion
of printed or handwritten text from images into machine-readable format, facilitating
document digitization, automated data entry, text search in scanned documents, and
accessibility features for visually impaired users.

Traditional \gls{ocr} systems relied heavily on hand-crafted features and rule-based
approaches, requiring extensive domain expertise and manual tuning for different fonts,
languages, and image conditions. The advent of deep learning has revolutionized this
field, enabling end-to-end learning systems that automatically discover optimal feature
representations directly from data.

\section*{Project Objectives}

This project aims to develop a complete \gls{ocr} system for printed character recognition
using modern deep learning techniques. The specific objectives are:

\begin{enumerate}
\item Design and implement a robust preprocessing pipeline to handle varied image
conditions including different lighting, contrast, and noise levels.

\item Build a deep neural network architecture combining convolutional and recurrent
layers to extract visual features and model sequential dependencies in text.

\item Train the model on a comprehensive dataset of 422,400 character images spanning
66 classes (digits 0-9, uppercase A-Z, lowercase a-z, and common punctuation marks).

\item Implement \gls{ctc} loss for sequence-to-sequence learning without requiring
explicit character-level alignment during training.

\item Deploy the trained model as a web application with a user-friendly interface
for real-time character and text recognition.

\item Evaluate model performance and identify areas for future improvement.
\end{enumerate}

\section*{Report Structure}

This report is organized as follows:

\textbf{Chapter 1} presents the problem statement, reviews related work in \gls{ocr}
and deep learning for text recognition, and discusses the theoretical foundations of
\gls{cnn}, \gls{rnn}, \gls{lstm}, and \gls{ctc} loss.

\textbf{Chapter 2} describes the dataset structure, preprocessing techniques including
binarization and normalization, and the data loading pipeline implemented for efficient
training.

\textbf{Chapter 3} details the model architecture, including the \gls{cnn} feature
extractor, \gls{bilstm} sequence encoder, and \gls{ctc} decoder. Implementation details
using TensorFlow and Keras are provided.

\textbf{Chapter 4} presents the training procedure, hyperparameter configuration,
deployment as a Flask web application, and evaluation results with discussion of
observed behaviors and limitations.

Finally, the \textbf{Conclusion} summarizes the achievements and proposes perspectives
for future enhancements.


% ──────────────────────────────────────────────────────────
%  CHAPITRES
% ──────────────────────────────────────────────────────────
\chapter{Introduction and Problem Statement}
\label{chap:chapterone}
\input{chapitre1.tex}

\chapter{Dataset and Preprocessing}
\label{chap:2}
\input{chapitre2.tex}

\chapter{Model Architecture and Implementation}
\label{chap:3}
\input{chapitre3.tex}

\chapter{Training, Deployment and Results}
\label{chap:4}
\input{chapitre4.tex}

% ──────────────────────────────────────────────────────────
%  CONCLUSION
% ──────────────────────────────────────────────────────────
\chapter*{Conclusion and Perspectives}
\label{chap:conclusion}
\addcontentsline{toc}{chapter}{Conclusion and Perspectives}
\markboth{\MakeUppercase{Conclusion and Perspectives}}{}

This project successfully developed an end-to-end \gls{ocr} system capable of recognizing
printed characters from images. We implemented a complete deep learning pipeline covering
data loading, preprocessing, model architecture design, training, and deployment as a
web application.

The proposed architecture — a \gls{crnn} combining a \gls{cnn} feature extractor with
stacked \gls{bilstm} layers and \gls{ctc} loss — demonstrated effectiveness for character
recognition tasks. Our dataset of 422,400 images across 66 character classes (digits,
uppercase, lowercase, and punctuation) provided sufficient training data for robust
character-level recognition.

The preprocessing pipeline, including Otsu binarization, aspect-ratio-preserving resizing,
and normalization, ensured consistent input quality. The Flask-based web interface with
real-time prediction capabilities demonstrates the practical applicability of the system.

Key technical achievements include:
\begin{itemize}
\item Successfully handling legacy \gls{hdf5} model weights in modern Keras 3 environment
\item Implementing automatic character segmentation via contour detection
\item Deploying a responsive web interface for real-time OCR predictions
\item Achieving functional character recognition with the trained model
\end{itemize}

\section*{Perspectives}

Future improvements that can enhance the system:

\begin{itemize}
\item \textbf{Architecture optimization:} Consider switching to a classification-based
approach for single-character recognition, which may provide better accuracy and faster
inference than sequence-based \gls{ctc} models for this specific use case.

\item \textbf{Data augmentation:} Expand training robustness through synthetic augmentation
including rotation, noise injection, perspective transforms, and elastic distortions to
better handle real-world image variations.

\item \textbf{Advanced decoding:} Integrate language model-based beam search decoding
to leverage contextual information and correct recognition errors in word-level predictions.

\item \textbf{Multi-line support:} Extend the system to handle full document pages by
integrating text detection models (CRAFT, DB-Net, or EAST) before the recognition stage.

\item \textbf{Model optimization:} Explore lightweight architectures for mobile deployment,
including MobileNet-based feature extractors or knowledge distillation techniques.

\item \textbf{Transformer architectures:} Investigate attention-based models such as
TrOCR or Vision Transformers for potentially improved recognition accuracy.

\item \textbf{Production deployment:} Implement model serving with TensorFlow Serving
or ONNX Runtime for scalable production deployment with proper monitoring and logging.
\end{itemize}

These perspectives open pathways for future work and the potential deployment of a
fully operational \gls{ocr} solution in real-world document digitization pipelines.

% ──────────────────────────────────────────────────────────
%  RÉFÉRENCES
% ──────────────────────────────────────────────────────────
\chapter*{References}
\addcontentsline{toc}{chapter}{References}

\begin{itemize}
\item Graves, A., Fernández, S., Gomez, F., \& Schmidhuber, J. (2006).
``Connectionist Temporal Classification: Labelling Unsegmented Sequence Data
with Recurrent Neural Networks.'' \textit{ICML 2006}.

\item Shi, B., Bai, X., \& Yao, C. (2016).
``An End-to-End Trainable Neural Network for Image-based Sequence Recognition
and Its Application to Scene Text Recognition.''
\textit{IEEE TPAMI}, 39(11), 2298--2304.

\item Abadi, M., et al.\ (2015).
\textit{TensorFlow: Large-Scale Machine Learning on Heterogeneous Systems}.
Software available from \texttt{tensorflow.org}.

\item Chollet, F., et al.\ (2015).
\textit{Keras}. Software available from \texttt{keras.io}.

\item Bradski, G. (2000).
``The OpenCV Library.'' \textit{Dr.\ Dobb's Journal of Software Tools}.

\item Otsu, N. (1979).
``A Threshold Selection Method from Gray-Level Histograms.''
\textit{IEEE Transactions on Systems, Man, and Cybernetics}, 9(1), 62--66.

\item LeCun, Y., Bengio, Y., \& Hinton, G. (2015).
``Deep Learning.'' \textit{Nature}, 521, 436--444.

\item Hochreiter, S., \& Schmidhuber, J. (1997).
``Long Short-Term Memory.'' \textit{Neural Computation}, 9(8), 1735--1780.

\item Pedregosa, F., et al.\ (2011).
``Scikit-learn: Machine Learning in Python.''
\textit{Journal of Machine Learning Research}, 12, 2825--2830.
\end{itemize}

\cleardoublepage
\printindex

\chapter*{Abstract}
\thispagestyle{empty}

\section*{English Abstract}

This project presents the design, implementation, and deployment of an Optical Character
Recognition (OCR) system for printed text using deep learning techniques. The system
employs a Convolutional Recurrent Neural Network (CRNN) architecture combining
Convolutional Neural Networks (CNN) for spatial feature extraction, Bidirectional
Long Short-Term Memory (BiLSTM) networks for temporal sequence modeling, and
Connectionist Temporal Classification (CTC) loss for end-to-end training without
explicit character alignment.

The dataset comprises 422,400 grayscale character images across 66 classes (digits,
uppercase and lowercase letters, and punctuation marks), with 6,400 images per class.
A comprehensive preprocessing pipeline including Otsu binarization, aspect-ratio-preserving
resizing, and normalization ensures robust input quality.

The trained model is deployed as a Flask web application providing real-time character
recognition through a user-friendly interface. The system implements automatic character
segmentation via contour detection for multi-character text recognition.

Key technical contributions include handling legacy HDF5 model formats in modern Keras 3
environments, implementing efficient data pipelines with tf.data, and deploying a
production-ready web interface with CORS support.

Evaluation reveals functional character recognition with identified areas for improvement,
including architecture optimization (switching from sequence models to classification
for single-character recognition) and data augmentation for enhanced robustness.

\keywordss{Optical Character Recognition, Deep Learning, Convolutional Neural Networks,
Recurrent Neural Networks, LSTM, CTC Loss, TensorFlow, Keras, Flask, Computer Vision}

\vspace{2cm}

\section*{Résumé en Français}

Ce projet présente la conception, l'implémentation et le déploiement d'un système de
Reconnaissance Optique de Caractères (OCR) pour texte imprimé utilisant des techniques
d'apprentissage profond. Le système emploie une architecture de Réseau de Neurones
Convolutif Récurrent (CRNN) combinant des Réseaux de Neurones Convolutifs (CNN) pour
l'extraction de caractéristiques spatiales, des réseaux Long Short-Term Memory
Bidirectionnels (BiLSTM) pour la modélisation de séquences temporelles, et la perte
Connectionist Temporal Classification (CTC) pour l'entraînement de bout en bout sans
alignement explicite des caractères.

Le jeu de données comprend 422 400 images de caractères en niveaux de gris réparties
sur 66 classes (chiffres, lettres majuscules et minuscules, et signes de ponctuation),
avec 6 400 images par classe. Un pipeline de prétraitement complet incluant la
binarisation d'Otsu, le redimensionnement préservant le ratio d'aspect et la
normalisation assure une qualité d'entrée robuste.

Le modèle entraîné est déployé comme application web Flask fournissant une reconnaissance
de caractères en temps réel via une interface conviviale. Le système implémente une
segmentation automatique des caractères via détection de contours pour la reconnaissance
de texte multi-caractères.

Les contributions techniques clés incluent la gestion des formats de modèles HDF5
hérités dans les environnements Keras 3 modernes, l'implémentation de pipelines de
données efficaces avec tf.data, et le déploiement d'une interface web prête pour la
production avec support CORS.

L'évaluation révèle une reconnaissance fonctionnelle des caractères avec des domaines
d'amélioration identifiés, notamment l'optimisation de l'architecture (passage des
modèles de séquence à la classification pour la reconnaissance de caractères uniques)
et l'augmentation de données pour une robustesse accrue.

\keywords{Reconnaissance Optique de Caractères, Apprentissage Profond, Réseaux de
Neurones Convolutifs, Réseaux de Neurones Récurrents, LSTM, Perte CTC, TensorFlow,
Keras, Flask, Vision par Ordinateur}

\end{document}


\end{document}
